\chapter{Tester og resultater}
\label{ch:results}
Dette kapitlet bevarer programvarebaselinen og tidligere observasjoner fra
2. november. Nye resultater fra 3. november står i
kapittel~\ref{ch:software-only}.

\section{Metode og evidensnivå}
Tester kjøres uten USB-/GPIO-tilgang og uten å programmere kortet i denne
gjennomgangen. Den faktiske Ada-koden kompileres nativt med assertions og
overflow-sjekker. Generiske grensesnitt og eksplisitte fakes erstatter bare
maskinvare og tid der dette er nødvendig. ARM-bygget kontrollerer at den
konkrete bindingen i \kode{Main} også kompilerer med riktige drivere.
Dette er komplementære kontroller; de må ikke telles som fysiske forsøk.

\begin{tabularx}{\textwidth}{@{}P{41mm}P{18mm}Y@{}}\toprule
Kontroll & Resultat & Omfang\\\midrule
Ada policy og grenser & 1729 & Eksisterende beslutninger, ugyldige data, overganger og grenseverdier.\\
Range-modell & 24 & Pulsmodell og gyldighetsgrenser.\\
Feilgrense & 38 & Primær feil, sekundær stoppfeil og loggerfeil.\\
Lyspolicy/status & 94 & Mønstre, sticky feil og beskyttet deling.\\
Adaptivt søk & 49 & Tidlig retur, geometri, alder og bevart stråletid.\\
Integrert kontroller & 521 & 22 scenarier med reell styringsflyt og injiserte callbacks.\\
Motor, simulert I2C & 76 & Alle 12 oppstartspunkter, arming og FULL\_OFF.\\
OLED, simulert I2C & 32 & Sunn oppstart, timeout og ingen nye forsøk etter feil.\\
Knapper & 6 & A/B, frigivelse og stoppforrang.\\
Navigasjon & 30 & Reelle sensor-/motorpakker mot fakes; stopp i avlesning og full tidsstyrt skanning.\\\midrule
Sum Ada & 2599 & Bestått; kontrollpunkter, ikke uavhengige kjøreforsøk.\\
Python & 40 & Analyse, protokoll- og loggverktøy.\\\bottomrule
\end{tabularx}

\section{Samspill og feilinjeksjon}
Kontrollertestene dekker venting uten A, stopp før arming, feil i
initialisering, gjentatt initialisering, fri Cruise uten stopp/start,
hindring, begge svingretninger, fri front etter skanning, blokkert område,
endret klarering, avvist kjørepermit og utløpt evidens før/under siste
kontroll. Videre testes B-stopp under sensorarbeid, taskfeil, motorfeil
og samtidig sekundær stoppfeil. Assertions krever at servo, tekst-I/O og
filtrert stasjonær kontroll ikke utføres mens testbilen er i bevegelse.

At en test passer er ikke nok til å vise at den oppdager feil. Fire
midlertidige kildekopier endres derfor bevisst: arming fjernes,
ferskhetskontroll fjernes, OLED-feilsperren fjernes og motorens FULL\_OFF
fjernes. Hver kopi må fortsatt kompilere og deretter bli avvist av riktig
testassertion. En femte mutasjon fjerner frontservoens ventefrist og
avvises av den reelle navigasjonstesten. Originalkoden endres ikke av
disse mutasjonene.

En full native skannerunde bruker reell 65 ms pacing og 500 ms servoventing
med en kontrollert sensorsekvens. Testen kontrollerer 24 råprøver,
sideavlesning under sentrering, ingen tidlig fronttrigger og at påfølgende
live-kontroll kan fullføres innen evidensvinduet på verten. Ubrukt tillatelse
utløper også. Vertstid er ikke målmaskinens WCET eller fysisk servotid.

\section{Avgrenset formell verifikasjon}
\kode{Adaptive_Lamp_Policy} har \kode{SPARK_Mode} i både spesifikasjon og
body. Funksjonen \kode{Outputs} har \kode{Global => null} og en
postcondition som blant annet krever rødt ved parkert/brems/feil,
sideblink bare i riktig modus og begge blinklys samtidig bare i Fault.
Det isolerte prosjektet analyserer akkurat denne produksjonsenheten
\cite{sparktutorial2026}.

GNATprove rapporterer fire samlede analysekontroller: dataavhengighet,
terminering, en kjøretidssjekk og en funksjonskontrakt, uten ubeviste
kontroller. En ekstra mutasjon som slår på bremselys i Moving må gi
avvist postcondition og ikke-null returkode. Beviskjøringen bruker
\kode{--checks-as-errors=on}. Beviset dekker ikke GPIO, taskplanlegging,
elektrisk strøm, motorstopp eller resten av navigasjonen.

\section{Bygg og reproduksjon}
Hovedprogram, sensor V2, GPIO V3 og lysdiagnostikk bygges separat med
GNAT ARM 15.1.2 og GPRbuild 25.0.1. Assertions og overflow-sjekker er
aktivert i egne firmwareprosjekter. GNATprove er låst til 15.1.0.
Bygg og tester gjentas med:
\begin{lstlisting}[language={}]
python3 tools/sync_report_sources.py --check
python3 tools/check_project.py --full
\end{lstlisting}
Kommandoen åpner ikke kortet. Loggene og et maskinlesbart manifest med
returkoder og kildehash lagres i \kode{build/readiness-review/}.
Vesentlige resultater kopieres til datert evidens i Git; genererte
firmwarefiler og PDF-er holdes utenfor versjonskontrollen.

\section{Tidligere fysiske observasjoner}
23. oktober ga tre stasjonære serier ved brukeroppgitt omtrent 60 cm totalt
480/480 gyldige venstremålinger. Middelverdien var 57,108 cm,
standardavvik omtrent 0,465 cm og avviket fra referansen $-2,892$ cm.
Dette viser repeterbarhet i det oppsettet, men også et avvik som ikke
kan korrigeres forsvarlig fra ett referansepunkt. Originaldata og analyser
er bevart i \kode{evidence/2026-10-23-reference60/}.

Tidligere 2. november ga GPIO V3 320 NO\_RISE-avlesninger mens
sensorforsyningen var av. Dette er en forventet observasjon av oppsettet,
ikke bevis for defekte sensorer. Senere lysdiagnostikk ga 362 godkjente
GPIO-registerprøver over seks mønstre. Ingen LED-er var montert.
Registerlatch er ikke en optisk eller elektrisk LED-test.

Disse prøvene gjelder de dokumenterte eldre firmwarehashene. De er ikke
gjentatt på dagens nye bygg. Kortet står fortsatt med den tidligere
lysdiagnostikken; den nye autonome versjonen er ikke lastet.

\section{Gjenstående fysisk akseptanse}
\begin{enumerate}
\item Monter og kontroller alle fire lys med egen motstand, riktig polaritet
og felles GND. Verifiser sidevalg og baklys visuelt og elektrisk.
\item Test oppstart med motorstrøm isolert. Lokaliser den åpne fysiske
I2C-oppstartsfeilen; kontroller både OLED tilkoblet og frakoblet.
\item Gi sensorene riktig forsyning. Kontroller ECHO-nivå og hver kanal ved
flere oppmålte avstander; registrer feilrate, bias, MAE og RMSE.
\item Løft hjulene og bekreft alle fortegn, begge svinger, A-arming og B-stopp.
\item Mål lavfartsbevegelse, bremsing og responstid før autonom løype.
Gjenta en sammenhengende sekvens med rålogg og video.
\end{enumerate}
Kriteriene og loggmalen ligger i \kode{KLAR-FOR-BENK.md}. Kravstatus i
kapittel~\ref{ch:requirements} oppdateres fra resultatene, ikke fra planen.
